% [原文 p.1]
\chapter{共形场理论的几何}

共形场理论是 1980 年代中期由 Belavin--Polyakov--Zamolodchikov 引入的场论，用来描述临界现象中出现的无限维对称性。本章解说共形场理论的几何侧面，特别是说明如何从 Wess--Zumino--Witten 模型出发，构造出共形变换下不变的场论。

此处起重要作用的是被称为环路群的无限维 Lie 群，以及其上的复线丛。首先，§1.1 解说环路群及其无穷小版本，即环路代数的中心扩张，并具体构造环路群上的复线丛。§1.2 汇总仿射 Lie 代数及其表示论，特别涉及 level $k$ 的最高权表示，以及由 Sugawara 形式给出的 Virasoro 代数表示。

§1.3 给出 Wess--Zumino--Witten 模型的表述。Wess--Zumino--Witten 模型对 Riemann 面 $\Sigma$ 到 Lie 群 $G$ 的光滑映射 $f$ 定义作用量 $S(f)$，并求配分函数，即 $\exp(-S(f))$ 关于所有 $f$ 的平均。这里，对带边界的 Riemann 面，用环路群上复线丛的语言来描述 $\exp(-S(f))$。§1.4 用仿射 Lie 代数的表示空间，把共形块空间定义为某种不变张量。共形场理论在曲面的模空间上构造出以共形块空间为纤维的向量丛及其上的联络，意义重大。这里特别在 §1.5 中，就 Riemann 球面的情形，说明点构形空间上的向量丛及其上称为 Knizhnik--Zamolodchikov 联络的平坦联络是如何导出的。

% [原文 p.2]
其次在 §1.6 中引入顶点算子的概念，并说明它与物理学家常用的算子积展开之间的关系。这里也涉及环面上的共形块空间、Verlinde 等式等。

% [Tutor Note: 本书的记号前缀不可混：$LG$ 是环路群（映射空间），$L\mathfrak{g}=\mathfrak{g}\otimes\mathbb{C}((t))$ 是与之对应的环路代数（形式 Laurent 版本），$\widehat{\mathfrak{g}}=L\mathfrak{g}\oplus\mathbb{C}c$ 是它的中心扩张。后文 $LG_{\mathbb{C}}$ 指 $G_{\mathbb{C}}$ 的环路群，与 $\widehat{\mathfrak{g}}$ 不是同一层次的对象。]

\section{环路群与仿射 Lie 代数}

\takeaway{本节从环路群 $LG$ 出发，造出环路代数 $L\mathfrak{g}$ 的中心扩张。核心工具是由 2 次上同调 $H^{2}(L\mathfrak{g},\mathbb{C})\cong\mathbb{C}$ 给出的 2-余循环；同一套余循环还给出仿射 Lie 代数、Virasoro 代数以及 $LG_{\mathbb{C}}$ 上的基本复线丛。}

首先定义本书中起重要作用的环路群。设 $G$ 为单连通紧单 Lie 群，把 $S^{1}=\{z\in\mathbb{C}\mid |z|=1\}$ 到 $G$ 的光滑映射全体记作 $LG$。在 $LG$ 上按
\[
  (\gamma_{1}\cdot\gamma_{2})(z)=\gamma_{1}(z)\gamma_{2}(z),\qquad \gamma_{1},\gamma_{2}\in LG
\]
引入积结构，所得对象称为\textbf{环路群}（loop group）。环路群 $LG$ 具有无限维 Lie 群的结构。以下，作为 Lie 群 $G$，主要处理 $SU(2)$ 的情形。

\textbf{环路代数}（loop algebra）指环路群 $LG$ 的 Lie 代数作复化所得的对象，但这里采用以 Laurent 级数写出的形式版本。所谓形式 Laurent 级数，指可用适当的 $m$ 写成 $f(t)=\sum_{n=-m}^{\infty}a_{n}t^{n}$ 者；把这类级数全体构成的 $\mathbb{C}$ 代数记作 $\mathbb{C}((t))$，并令
\[
  L\mathfrak{g}=\mathfrak{g}\otimes\mathbb{C}((t)).
\]
这里 $\mathfrak{g}$ 是 $G$ 的 Lie 代数的复化。$t$ 是形式变量，今后有时也把它看作 $\mathbb{C}\setminus\{0\}$ 的参数。

对 $L\mathfrak{g}$ 的元素 $X\otimes f$、$Y\otimes g$，按
\begin{equation}
  [X\otimes f,\,Y\otimes g]=[X,Y]\otimes fg \tag{1.1}
\end{equation}
定义 Lie 括号，$L\mathfrak{g}$ 便具有 Lie 代数的结构。

下面说明 $L\mathfrak{g}$ 的中心扩张。作为 $\mathbb{C}$ 上的线性空间，把 $\widehat{\mathfrak{g}}$ 定为 $L\mathfrak{g}$ 与以 $c$ 为基的 $\mathbb{C}$ 上一维线性空间的直和 $L\mathfrak{g}\oplus\mathbb{C}c$。以下在 $\widehat{\mathfrak{g}}=L\mathfrak{g}\oplus\mathbb{C}c$ 上引入使 $c$ 属于中心的 Lie 代数结构。这里，$c$ 属于中心是指对一切 $\xi\in\widehat{\mathfrak{g}}$ 都有 $[c,\xi]=0$。

% [原文 p.3]
利用双线性映射 $\omega\colon L\mathfrak{g}\times L\mathfrak{g}\to\mathbb{C}$，把 $\widehat{\mathfrak{g}}$ 的括号定为
\[
  [\xi+\alpha c,\,\eta+\beta c]=[\xi,\eta]+\omega(\xi,\eta)c,
  \qquad \xi,\eta\in L\mathfrak{g},\ \alpha,\beta\in\mathbb{C}
\]
时，它真正给出 Lie 代数结构的条件是：对一切 $x,y,z\in L\mathfrak{g}$，
\begin{equation}
  \omega(x,y)=-\omega(y,x) \tag{1.2}
\end{equation}
\begin{equation}
  \omega([x,y],z)+\omega([y,z],x)+\omega([z,x],y)=0 \tag{1.3}
\end{equation}
均被满足。事实上，条件 (1.2) 对应括号的反对称性 $[\xi,\eta]=-[\eta,\xi]$，条件 (1.3) 则对应 Jacobi 等式。

对这样的 $\omega$，$\widehat{\mathfrak{g}}$ 的 Lie 括号由
\[
  [X\otimes f,\,Y\otimes g]=[X,Y]\otimes fg+\omega(X\otimes f,\,Y\otimes g)c
\]
给出。关于 $\omega$ 的条件 (1.3) 称为 \textbf{2-余循环条件}（2-cocycle condition）。又，满足条件 (1.2)、(1.3) 的双线性映射 $\omega$、$\omega'$ 在 $L\mathfrak{g}\oplus\mathbb{C}c$ 上定出同构的 Lie 代数结构的充要条件是：存在线性映射 $\mu\colon L\mathfrak{g}\to\mathbb{C}$，使对一切 $x,y\in L\mathfrak{g}$ 有
\[
  \omega(x,y)=\omega'(x,y)+\mu([x,y]).
\]
这样一来，对满足条件 (1.2)、(1.3) 的 $\omega$，可以在 $L\mathfrak{g}\oplus\mathbb{C}c$ 上引入使 $c$ 属于中心的 Lie 代数结构；这个结构称为 $L\mathfrak{g}$ 的\textbf{中心扩张}（central extension）。

一般地，先准备好 Lie 代数上同调的语言。对复 Lie 代数 $\mathfrak{a}$ 与左 $\mathfrak{a}$-模 $M$，把 $\mathfrak{a}$ 的以 $M$ 为系数的 $p$ 次上链群 $C^{p}(\mathfrak{a},M)$ 定义为 $\mathrm{Hom}_{\mathbb{C}}\bigl(\bigwedge^{p}\mathfrak{a},M\bigr)$。微分 $d_{p}\colon C^{p}(\mathfrak{a},M)\to C^{p+1}(\mathfrak{a},M)$ 对 $\omega\in C^{p}(\mathfrak{a},M)$ 由
\[
\begin{gathered}
  (d\omega)(x_{0},x_{1},\dots,x_{p}) \\
  =\sum_{i=0}^{p}(-1)^{i}x_{i}\omega(x_{0},\dots,\widehat{x_{i}},\dots,x_{p}) \\
  \qquad +\sum_{0\le i<j\le p}(-1)^{i+j}
  \omega([x_{i},x_{j}],x_{0},\dots,\widehat{x_{i}},\dots,\widehat{x_{j}},\dots,x_{p})
\end{gathered}
\]
确定。这里 $\widehat{x_{i}}$ 表示去掉分量 $x_{i}$。

% [原文 p.4]
如此得到的复形的上同调群
\[
  H^{p}(\mathfrak{a},M)=\mathrm{Ker}\,d_{p}/\mathrm{Im}\,d_{p-1}
\]
称为 $\mathfrak{a}$ 的以 $M$ 为系数的 $p$ 次上同调群。

如上所述，$H^{2}(L\mathfrak{g},\mathbb{C})$ 的元素定出环路代数 $L\mathfrak{g}$ 的中心扩张的同构类。反之，设
\[
  0\longrightarrow\mathbb{C}\longrightarrow\widetilde{L\mathfrak{g}}
  \xrightarrow{\ \pi\ }L\mathfrak{g}\longrightarrow 0
\]
为 $L\mathfrak{g}$ 的中心扩张。取线性空间意义下的截面 $s\colon L\mathfrak{g}\to\widetilde{L\mathfrak{g}}$，即满足 $\pi\circ s=\mathrm{id}$ 的线性映射。对 $Z\in L\mathfrak{g}$ 记 $s(Z)=\widetilde{Z}$，并令
\[
  \eta(X,Y)=[\widetilde{X},\widetilde{Y}]-\widetilde{[X,Y]},
  \qquad X,Y\in L\mathfrak{g},
\]
则 $\eta$ 与 $s$ 的取法无关，定出 $H^{2}(L\mathfrak{g},\mathbb{C})$ 的一个元素。

这样一来，环路代数 $L\mathfrak{g}$ 的中心扩张的同构类与 $H^{2}(L\mathfrak{g},\mathbb{C})$ 的元素之间存在一一对应；对一般的 Lie 代数这也是熟知的事实。把 Lie 代数 $\mathfrak{g}$ 的 Cartan--Killing 形式记作
\[
  \langle\ ,\ \rangle\colon \mathfrak{g}\times\mathfrak{g}\longrightarrow\mathbb{C}.
\]
它是对称双线性形式，且在 $\mathfrak{g}$ 的伴随作用下不变，即对任意 $X,Y,Z\in\mathfrak{g}$ 有
\[
  \langle [X,Y],Z\rangle=\langle X,[Y,Z]\rangle .
\]
例如 §1.2 也将说明，当 $\mathfrak{g}=\mathfrak{sl}_{2}(\mathbb{C})$ 时 $\langle X,Y\rangle=\mathrm{Tr}(XY)$ 定出 $\mathfrak{g}$ 的 Cartan--Killing 形式。对环路代数 $L\mathfrak{g}$，确实可以如下具体构造 2-余循环。

\begin{proposition*}{1.1}
  对环路代数 $L\mathfrak{g}$，有 $H^{2}(L\mathfrak{g},\mathbb{C})\cong\mathbb{C}$，且它由
  \[
    \omega(X\otimes f,\,Y\otimes g)=\langle X,Y\rangle\,
    \mathrm{Res}_{t=0}(df\,g)
  \]
  定出的 2-余循环 $\omega$ 生成。这里取
  $\mathrm{Res}_{t=0}\bigl(\sum_{n}c_{n}t^{n}dt\bigr)=c_{-1}$。
\end{proposition*}

\begin{proof}
  首先证明：对 2-余循环 $\alpha$，$\alpha$ 所属的上同调类可由在 Lie 群 $G$ 的共轭作用下不变的元来代表。把 $G$ 的元素 $g$ 写成 $\exp(tZ)$，$Z\in\mathfrak{g}$，则对 $X\in\mathfrak{g}$ 有
  $gXg^{-1}=X+t[Z,X]+O(t^{2})$，于是用等式 (1.3) 可得

% [原文 p.5]
  \[
    \lim_{t\to 0}\frac{1}{t}\bigl[\alpha(gXg^{-1},gYg^{-1})-\alpha(X,Y)\bigr]
    =\alpha(Z,[X,Y]).
  \]
  此式右端在令 $\mu_{Z}\colon L\mathfrak{g}\to\mathbb{C}$、$\mu_{Z}(U)=\alpha(Z,U)$ 时可写成 $\mu_{Z}([X,Y])$。因此，经 Lie 群 $G$ 的共轭作用的无穷小版本，上同调类不改变。利用这一点，用 $G$ 的不变测度对上述作用取平均，便知可唯一地得到 $G$ 不变的上同调代表元。

  对 2-余循环 $\alpha\colon\bigwedge^{2}L\mathfrak{g}\to\mathbb{C}$，令
  $\alpha_{m,n}(X,Y)=\alpha(X\otimes t^{m},Y\otimes t^{n})$。这里不妨设 $\alpha_{m,n}$ 是 Lie 群 $G$ 的 Lie 代数复化 $\mathfrak{g}$ 上 $\mathfrak{g}$ 不变的双线性形式。由 $\mathfrak{g}$ 是单 Lie 代数可知它是对称形式。又由 $\alpha$ 是反对称形式可知 $\alpha_{m,n}=-\alpha_{n,m}$。此外，$\alpha$ 为 2-余循环这一条件表为
  \[
    \alpha_{m+n,p}+\alpha_{n+p,m}+\alpha_{p+m,n}=0.
  \]
  先令 $n=p=0$，得 $\alpha_{m,0}=0$。取 $p=-m-n$、$n=1$ 时，由上式得到的递推式可导出 $\alpha_{m,-m}=m\alpha_{1,-1}$。再令 $p=q-m-n$，则
  $\alpha_{q-m-n,m+n}=\alpha_{q-m,m}+\alpha_{q-n,n}$，用它得到
  $\alpha_{q-k,k}=k\alpha_{q-1,1}$。进而由 $q\alpha_{q-1,1}=\alpha_{0,q}=0$ 推出：当 $m+n\ne 0$ 时 $\alpha_{m,n}=0$。综上可知
  \[
    \alpha_{m,n}=m\delta_{m+n,0}\alpha_{1,-1}.
  \]
  $\alpha_{1,-1}\colon\mathfrak{g}\times\mathfrak{g}\to\mathbb{C}$ 是 $\mathfrak{g}$ 不变的内积，故除常数倍外只能是 Cartan--Killing 形式。取 $\alpha_{1,-1}$ 为 Cartan--Killing 形式，即得所求 2-余循环 $\omega$ 的形状。这个 2-余循环作为上同调类非平凡，例如可由下述看出：取 $\mathfrak{g}$ 的 Cartan 子代数的一个基元 $H$，对环路代数的 Lie 括号有 $[H\otimes t,H\otimes t^{-1}]=0$，而 $\langle H,H\rangle\ne 0$，由此立即得出该上同调类非平凡。以上证明了命题。

% [Tutor Note: 末句用 $[H\otimes t,H\otimes t^{-1}]=0$ 而余循环值 $\langle H,H\rangle\ne 0$，说明该扩张不能靠 $\mu([x,y])$ 消去，即上同调类非平凡；关键是括号为零而中心项不为零。]
\end{proof}

% [Tutor Note: 证明中的 $t$ 是 $\exp(tZ)$ 的参数，与形式变量 $t$ 无关，只是符号被重复使用；读 $gXg^{-1}=X+t[Z,X]+O(t^{2})$ 时不要把它当作 Laurent 变量。]

% [Tutor Note: 命题 1.1 中 $df\,g$ 的 $d$ 只作用于 $f$；若写成 $d(fg)$ 或 $f\,dg$，因全微分的留数为零会差一个符号，$\mathrm{Res}(df\,g)=-\mathrm{Res}(f\,dg)$。]

把 $\widehat{\mathfrak{g}}$ 称为伴随于 Lie 代数 $\mathfrak{g}$ 的\textbf{仿射 Lie 代数}（affine Lie algebra）。这里定下的 Lie 括号也可写成如下形状：

\begin{equation}
  [X\otimes t^{m},\,Y\otimes t^{n}]=[X,Y]\otimes t^{m+n}
  +\langle X,Y\rangle\,m\delta_{m+n,0}\,c \tag{1.4}
\end{equation}

% [原文 p.6]
下面考虑从复平面的单位圆 $S^{1}$ 到自身、保持定向的微分同胚全体，记作 $\Diff^{+}(S^{1})$。$\Diff^{+}(S^{1})$ 具有无限维 Lie 群的结构，其 Lie 代数是 $S^{1}$ 上光滑向量场全体 $\mathrm{Vect}\,S^{1}$。在复化 $[\mathrm{Vect}\,S^{1}]\otimes\mathbb{C}$ 的元素中，考虑那些可用 Laurent 多项式 $f(z)$ 写成
\[
  f(z)\frac{d}{dz}
\]
的形状、并能延拓到 $\mathbb{C}\setminus\{0\}$ 的向量场全体，令
\[
  \cA=\Bigl\{\,f(z)\frac{d}{dz}\ \Bigm|\ f(z)\in\mathbb{C}[z,z^{-1}]\,\Bigr\}.
\]
这里 $\mathbb{C}[z,z^{-1}]$ 表示 Laurent 多项式全体构成的 $\mathbb{C}$ 代数。$\cA$ 按
\[
  \Bigl[f(z)\frac{d}{dz},\,g(z)\frac{d}{dz}\Bigr]
  =\{f(z)g'(z)-g(z)f'(z)\}\frac{d}{dz}
\]
定义 Lie 括号，便具有 Lie 代数的结构。这是无限维 Lie 代数，以
\[
  L_{n}=-z^{n+1}\frac{d}{dz},\qquad n\in\Z
\]
为基。向量场 $L_{n}$ 令 $z=e^{i\theta}$ 并限制到 $S^{1}$ 上，可表为
\[
  L_{n}=ie^{in\theta}\frac{\partial}{\partial\theta}.
\]
关于上面的 Lie 括号，$L_{n}$（$n\in\Z$）满足
\begin{equation}
  [L_{m},L_{n}]=(m-n)L_{m+n} \tag{1.5}
\end{equation}
此外，向量场 $L_{n}$ 生成 $\mathbb{C}\setminus\{0\}$ 的无穷小变换
\[
  \varphi_{t}(z)=z-tz^{n+1},
\]
而在 $L_{n}$（$n\in\Z$）中，能作为 Riemann 球面 $\CP^{1}$ 整体的光滑向量场延拓的只有 $L_{0}$、$L_{-1}$、$L_{1}$。它们满足
\[
  [L_{0},L_{1}]=-L_{1},\qquad [L_{0},L_{-1}]=L_{-1},\qquad
  [L_{1},L_{-1}]=2L_{0},
\]
并生成 Lie 代数 $\mathfrak{sl}_{2}(\mathbb{C})$。与仿射 Lie 代数的情形相同，用满足 2-余循环条件的反对称双线性映射 $\alpha\colon\cA\times\cA\to\mathbb{C}$ 可以构造 $\cA$ 的中心扩张。

% [原文 p.7]
下述命题成立。

\begin{proposition*}{1.2}
  对 Laurent 多项式向量场构成的 Lie 代数 $\cA$，有 $H^{2}(\cA,\mathbb{C})\cong\mathbb{C}$，且它由
  \[
    \omega\Bigl(f\frac{d}{dz},\,g\frac{d}{dz}\Bigr)
    =\frac{1}{12}\,\mathrm{Res}_{z=0}(f'''g\,dz)
  \]
  定出的 2-余循环 $\omega$ 生成。
\end{proposition*}

\begin{proof}
  取 2-余循环 $\alpha$，令 $\alpha_{p,q}=\alpha(L_{p},L_{q})$。对 $(L_{0},L_{p},L_{q})$ 写出 $\alpha$ 的 2-余循环条件，便知 $\alpha$ 的上同调类在以原点为中心的旋转下不变。因此取平均后，不妨设 $\alpha$ 自身在旋转下不变。此时若 $p+q\ne 0$ 则 $\alpha_{p,q}=0$，故令 $\alpha_{p}=\alpha_{p,-p}$。利用 $\alpha_{-p}=-\alpha_{p}$，2-余循环条件表为
  \[
    (p+2q)\alpha_{p}-(2p+q)\alpha_{q}=(p-q)\alpha_{p+q}.
  \]
  由此可知 $\alpha_{p}$ 具有 $\alpha_{p}=\lambda p^{3}+\mu p$ 的形状。改变 $\mu$ 的值不改变上同调类；特别取 $\lambda=\frac{1}{12}$、$\mu=-\frac{1}{12}$，即得所求的 2-余循环 $\omega$。
\end{proof}

用命题 1.2 的 2-余循环 $\omega$，在 $V=\cA\oplus\mathbb{C}c$ 上按
\[
  \Bigl[f\frac{d}{dz}+\alpha c,\ g\frac{d}{dz}+\beta c\Bigr]
  =\Bigl[f\frac{d}{dz},\,g\frac{d}{dz}\Bigr]
  +\omega\Bigl(f\frac{d}{dz},\,g\frac{d}{dz}\Bigr)c,
  \qquad \alpha,\beta\in\mathbb{C}
\]
定义 Lie 括号，便得到 $\cA$ 的中心扩张。称 $V$ 为 \textbf{Virasoro 代数}（Virasoro algebra）。Virasoro 代数的 Lie 括号可表为
\begin{equation}
  [L_{m},L_{n}]=(m-n)L_{m+n}+\frac{m^{3}-m}{12}\,\delta_{m+n,0}\,c. \tag{1.6}
\end{equation}

% [Tutor Note: 式 (1.6) 中中心元 $c$ 的系数是 $(m^{3}-m)/12$，即把中心荷归一到 1；文献里常见写法是 $(c/12)(m^{3}-m)$ 或 $(c/12)(m^{3}-m)\delta_{m+n,0}$，与此处「$c$ 本身是基元」的约定差一个倍数，换算时勿混。]

下面说明 Lie 群 $G$ 的复化 $G_{\C}$ 的环路群 $LG_{\C}$ 上的复线丛，以及 $LG_{\C}$ 的中心扩张。这些在后面的 Wess--Zumino--Witten 模型中起基本作用。一般地，取 $M$ 为微分流形，$L$ 为 $M$ 上的复线丛，并固定 $L$ 的联络 $\nabla$。

% [原文 p.8]
局部地，设 $M$ 的开集 $U$ 上复线丛 $L$ 已与 $U\times\mathbb{C}$ 平凡化，在 $U$ 上把联络 $\nabla$ 写成
\[
  \nabla=d-2\pi\sqrt{-1}\,\alpha .
\]
这里 $\alpha$ 是 $U$ 上定义的 1 次微分形式。用使 $L$ 局部平凡的开集覆盖 $M$，并在每个开集上作上述构造，则 $d\alpha$ 便在 $M$ 上定出 2 次闭微分形式。称它为联络 $\nabla$ 的\textbf{第一 Chern 形式}，记作 $c_{1}(\nabla)$。把 $c_{1}(\nabla)$ 写作 $\omega$，则众所周知，$\omega$ 满足 $d\omega=0$，且 $\omega$ 定出的上同调类不依赖联络 $\nabla$ 的选取，而定出 $H^{2}(M,\Z)$ 的一个元素。这就是 $L$ 的第一 Chern 类。

设 $M$ 为单连通微分流形。以 $M$ 上的点 $x_{0}$ 为基点、沿 $M$ 的环路对联络 $\nabla$ 所作的和乐，可如下描述。设 $D$ 是以 $\gamma$ 为边界、带定向的二维圆盘，则从 $L$ 在 $x_{0}$ 上的纤维 $L_{x_{0}}$ 的元素 $u$ 出发，沿 $\gamma$ 作关于联络 $\nabla$ 的平行移动，再回到 $x_{0}$ 所得的结果，由 Stokes 定理表为
\[
  u\exp\Bigl(2\pi\sqrt{-1}\int_{D}c_{1}(\nabla)\Bigr).
\]

% [Tutor Note: 式中的 $2\pi\sqrt{-1}$ 不能省成 $\sqrt{-1}$：正是这个因子使 $d\alpha$ 的 de Rham 类落在 $H^{2}(M,\Z)$ 的像里，即整性条件与 $2\pi\sqrt{-1}$ 的归一化绑定。]

作为逆的构造，我们来证明下述命题。

\begin{proposition*}{1.3}
  设 $M$ 为单连通微分流形。取 $M$ 上的 2 次微分形式 $\omega$，使 $d\omega=0$，且 $\omega$ 的上同调类定出 $H^{2}(M,\Z)$ 的元素。也就是说，设 $\omega$ 的 de Rham 上同调类属于自然映射
  \[
    i\colon H^{2}(M,\Z)\longrightarrow H^{2}(M,\R)
  \]
  的像。此时，$M$ 上的复线丛 $L$ 及其联络 $\nabla$ 中，第一 Chern 形式 $c_{1}(\nabla)$ 与 $\omega$ 一致者被唯一确定。
\end{proposition*}

\begin{proof}
  如下构造复线丛 $L$。把以 $x_{0}$ 为起点的 $M$ 的光滑道路全体记作 $P_{x_{0}}(M)$，即 $P_{x_{0}}(M)$ 是满足 $p(0)=x_{0}$ 的光滑映射 $p\colon[0,1]\to M$ 全体。在直积 $P_{x_{0}}(M)\times\mathbb{C}$ 上如下引入等价关系 $\sim$：对 $P_{x_{0}}(M)\times\mathbb{C}$ 的元素 $(p,u)$、$(q,v)$，规定 $(p,u)\sim(q,v)$ 是指 $p(1)=q(1)$，并且

% [原文 p.9]
  \[
    v=u\exp\Bigl(2\pi\sqrt{-1}\int_{D}\omega\Bigr)
  \]
  成立。这里 $D$ 是 $M$ 的带定向二维圆盘，其边界为 $p$ 与 $q^{-1}$ 复合所得的环路。由 $\omega$ 的上同调类定出 $H^{2}(M,\Z)$ 的元素这一假设，上述定义不依赖 $D$ 的选取而有意义。令 $L=P_{x_{0}}(M)\times\mathbb{C}/\sim$，并由 $\pi(p,u)=p(1)$ 定义投影 $\pi\colon L\to M$，则 $L$ 具有 $M$ 上复线丛的结构。由等价关系的定义，$L$ 上指定了平行移动，因而自然定出联络；又由上面的考察，这个联络的第一 Chern 形式与 $\omega$ 一致。
\end{proof}

以下取 $G=SU(2)$，描述环路群 $LG_{\C}$ 上的复线丛。$G_{\C}$ 是 $G$ 的复化，此时 $G_{\C}=SL(2,\mathbb{C})$。先证明环路群 $LG$ 单连通。把满足 $\gamma(1)=e$ 的 $\gamma\in LG$ 全体记作 $\Omega G$。又把 $G$ 作为常值映射嵌入 $LG$。由于 $LG$ 与直积 $G\times\Omega G$ 同胚，故
\[
  \pi_{1}(LG)\cong\pi_{1}(G)\oplus\pi_{1}(\Omega G)\cong\pi_{1}(G)\oplus\pi_{2}(G)
\]
成立。这里 $G$ 单连通且 $\pi_{2}(G)=0$，故知 $LG$ 单连通。同法可证 $LG_{\C}$ 单连通。于是把命题 1.3 的构造施于 $LG_{\C}$：对 $LG_{\C}$ 上满足 $d\omega=0$、且上同调类定出 $H^{2}(LG_{\C},\Z)$ 元素的 2 次微分形式 $\omega$，可构造 $LG_{\C}$ 上的复线丛。

下面给出这种 2 次微分形式的具体形状。$LG_{\C}$ 的 Lie 代数 $\Lie(LG_{\C})$ 可看作 $S^{1}$ 到 $\mathfrak{g}$ 的光滑映射全体，其中 $\mathfrak{g}$ 是 $G_{\C}$ 的 Lie 代数 $\mathfrak{sl}_{2}(\mathbb{C})$。把这样的映射作 Fourier 展开，并按与命题 1.1 相同的方法，对 $\xi,\eta\in\Lie(LG_{\C})$ 由
\begin{equation}
  \omega(\xi,\eta)=\frac{1}{4\pi^{2}}\int_{0}^{2\pi}
  \langle \xi'(\theta),\eta(\theta)\rangle\,d\theta \tag{1.7}
\end{equation}
定出 $\Lie(LG_{\C})$ 的 2-余循环。容易验证它确实给出 Lie 代数 $\Lie(LG_{\C})$ 的 2-余循环。此外，$\omega$ 定出 Lie 群 $LG_{\C}$ 上的左不变微分形式，且由 2-余循环条件知它是闭微分形式。

% [Tutor Note: $\pi_{1}(\Omega G)\cong\pi_{2}(G)$ 来自道路空间纤维化 $PG\to G$ 的同伦长正合列；这是「$LG$ 单连通」的全部依据，书中未展开。]

% [原文 p.10]
这个微分形式 $\omega$ 给出整系数上同调 $H^{2}(LG_{\C},\Z)$ 的元素，可如下说明。对 $SU(2)$ 的元素 $X$，$X^{-1}dX$ 是取值于 Lie 代数 $\mathfrak{su}(2)$ 的 $SU(2)$ 上 1 次微分形式，且关于 $SU(2)$ 的左作用不变。微分形式 $X^{-1}dX$ 称为 $SU(2)$ 的 Maurer--Cartan 形式。三次微分形式
\begin{equation}
  \sigma=\frac{1}{24\pi^{2}}\Tr\bigl(X^{-1}dX\wedge X^{-1}dX\wedge X^{-1}dX\bigr) \tag{1.8}
\end{equation}
定出的 de Rham 上同调类是 $H^{3}(SU(2),\Z)$ 的生成元。

% [Tutor Note: 式 (1.8) 的 $1/(24\pi^{2})$ 是使 $\sigma$ 在 $SU(2)\cong S^{3}$ 上积分为 1 的归一化，从而 $\sigma$ 的 de Rham 类是 $H^{3}(SU(2),\Z)$ 的整生成元；常数若换成 $1/8\pi^{2}$ 等便失去整性。]

如下定义 $LG$ 上的微分形式 $\omega_{0}$：由 $\phi(\gamma,z)=\gamma(z)$ 定义映射 $\phi\colon LG\times S^{1}\to G$，并令
\[
  \omega_{0}=-\int_{S^{1}}\phi^{*}\sigma .
\]
由定义，$\omega_{0}$ 满足 $d\omega_{0}=0$，其 de Rham 上同调类定出 $H^{2}(LG,\Z)$ 的元素。又 $\omega_{0}$ 自然延拓到 $LG_{\C}$。

把 $LG_{\C}$ 在 $\gamma\in LG_{\C}$ 处的切空间记作 $T_{\gamma}LG_{\C}$，并借助左不变向量场把它与 $\Lie(LG_{\C})$ 等同。把 $LG_{\C}$ 上的一次微分形式 $\beta$ 定义为：对 $\gamma\in LG_{\C}$、$\xi\in T_{\gamma}LG_{\C}$，令
\[
  \beta_{\gamma}(\xi)=\frac{1}{8\pi^{2}}\int_{0}^{2\pi}
  \bigl\langle \gamma(\theta)^{-1}\gamma'(\theta),\,
  \gamma(\theta)^{-1}\xi(\theta)\bigr\rangle\,d\theta .
\]
此时可以验证 $\omega=\omega_{0}+d\beta$ 成立。因此 $\omega$ 也定出整系数上同调类。综合以上讨论可知：可以构造 $LG_{\C}$ 上的复线丛及其联络，使其第一 Chern 形式为 $\omega$。把这个复线丛称为 $LG_{\C}$ 上的\textbf{基本复线丛}。事实上，已知 $H^{2}(LG,\Z)\cong\Z$，且它由式 (1.7) 所定义的 $\omega$ 的 de Rham 上同调类生成。

% [Tutor Note: $LG$ 单连通只保证 $H^{1}(LG,\Z)=0$，$H^{2}$ 一般不为零。事实上 $LG\cong G\times\Omega G$，故 $H^{2}(LG,\Z)\cong H^{2}(G,\Z)\oplus H^{2}(\Omega G,\Z)\cong 0\oplus\Hom(\pi_{3}(G),\Z)\cong\Z$；同理 $H^{2}(LG_{\C},\Z)\cong\Z$。末句把 (1.7) 的 $\omega$ 与其上同调类等同看待（本节已证 $\omega=\omega_{0}+d\beta$，$\omega_{0}$ 的类属于 $H^{2}(LG,\Z)$），并非笔误。]

% CHUNK ch01a DONE
